\documentclass[sigconf,nonacm]{acmart}
\settopmatter{printacmref=false,printfolios=true}
\usepackage{xurl}
\setlength{\emergencystretch}{2em}
\title{Benchmarking Freeway Estimation Under Structured Detector Loss}
\subtitle{ACM KDD 2027 | Datasets and Benchmarks working draft}
\author{Arjun Vijay Prakash}
\affiliation{\institution{Independent Researcher}\city{ }\country{ }}

\begin{document}
\begin{abstract}
Point error on randomly hidden readings can miss failures that arise during a contiguous
detector outage. CREST exports masked speed, flow, and storage and an evaluator-side scorer.
The benchmark has 33 synthetic scenarios split by seed and nine measured cases derived from
the three available FHWA I-80 and US-101 recordings. In a synthetic comparison with 360 targets
per mask regime, persistence outperformed the tree on random masks and the tree outperformed
persistence on contiguous masks. The benchmark reports coverage, fallback use, and interval
width alongside point error.
\end{abstract}

\maketitle

\section*{Tasks and split construction}
The simulator generates a time grid, directed incidence, capacity, and known traffic state. The
exporter constructs three contiguous outage blocks per seed at different positions along the
corridor and at different times. The exported splits contain 12 training scenarios across four
seeds, 9 calibration scenarios across three seeds, and 12 evaluation scenarios across four
seeds. Developing on chains and evaluating on junctions tests sensitivity to a change in network
layout. The cases from a seed share one truth realization, so uncertainty analysis must group by
seed.

Participant files contain observations, visibility, time, cell positions, and graph incidence.
Capacity is provided for synthetic cases. Evaluator files contain targets and the eligibility
mask. A checksummed manifest identifies each seed or recording cluster. A blind evaluation
requires evaluator-controlled seeds or recordings because the supplied sources are public.

\section*{Scoring and an illustrative ranking change}
Predictions have the same channel, time, and cell axes as the inputs. The scorer reports mean
absolute error for emitted points, eligible and scored counts, abstentions, and the number of
scored fallbacks. Paired lower and upper estimates add coverage, interval counts, and mean
width. NaN represents a withheld output. Split-specific target files stay on the evaluator side
so a participant input can be inspected directly for leakage.

In a separate matched experiment, every contiguous mask had a random control of equal size.
Across four synthetic seed clusters and 360 targets per regime, speed MAE was 1.63 km/h for
persistence and 7.18 for the tree under random masks. Under contiguous loss it was 8.21 and
7.28 km/h, respectively. The export scorer also ran historical median, persistence, and
interpolation, plus an observed-only control to verify the abstention denominator.

\section*{External validity and access}
The synthetic fixture samples every 30 seconds. The measured export uses one-minute aggregates
of FHWA trajectories. The two I-80 recordings are assigned to training and calibration, and
US-101 is the evaluation recording. Median observed speeds were 80.03 km/h in the
synthetic profile and 28.32 km/h across the three measured recordings. The measured bundle
carries FHWA attribution and CC BY-SA 3.0 terms. PeMS serves as a local comparison source.

\section*{Reproduction}
Run \texttt{python -m src.research.roadmap export-benchmark} for the synthetic bundle.
Add \texttt{--source ngsim} to export the trajectory bundle, or run \texttt{ablate} for
the matched-mask study.

\section*{References}
[1] CREST implementation guide and benchmark exchange format.
\url{https://github.com/ArjunCodess/CREST/blob/main/docs/crest-implementation.md}

[2] FHWA. Next Generation Simulation Vehicle Trajectories and Supporting Data.
\url{https://data.transportation.gov/d/8ect-6jqj}
\end{document}
